\chapter{Guía de estudio (Categoría 0)}

\begin{introcategoria}
Categoría obligatoria. Debe incluir, \textbf{para cada tema}, un calendario con el plan
docente detallado de \textbf{cada sesión}: materiales utilizados, contenidos cubiertos,
objetivos de aprendizaje, actividades planificadas y recomendación de trabajo personal.
\emph{No es la guía docente oficial de la asignatura.}
\end{introcategoria}

% Una tabla por tema directamente en LaTeX (ejemplo) ...
\section{Plan de sesiones}

\noindent\textbf{Tema 1. Título del tema}\\[4pt]
\begin{tabularx}{\textwidth}{@{}l X X@{}}
\toprule
\textbf{Sesión} & \textbf{Contenidos y objetivos} & \textbf{Materiales, actividades y trabajo personal}\\
\midrule
S1 (semana 1) & Contenido \ldots \newline \emph{Objetivo:} \ldots & Presentación T1 (cap.~3), apuntes \S1.1 \newline Actividad: \ldots \newline Trabajo personal: \ldots\\
S2 (semana 1) & \ldots & \ldots\\
\bottomrule
\end{tabularx}

% ... o incluye tus PDFs ya hechos:
\material{Presentación de la asignatura}{materiales/0-guia/presentacion.pdf}
  {Visión general de la asignatura: objetivos, contenidos, metodología y evaluación.}
\material{Planificación del curso}{materiales/0-guia/planificacion.pdf}
  {Calendario orientativo con la distribución de temas, sesiones y actividades.}
